% ex: ts=2 sw=2 sts=2 et filetype=tex
% SPDX-License-Identifier: CC-BY-SA-4.0

\documentclass[aspectratio=169]{beamer}

\usetheme{Boadilla}

\definecolor{ccosscolor}{rgb}{0.0, 0.545, 0.545}
\usecolortheme[named=ccosscolor]{structure}

\setbeamertemplate{navigation symbols}{} % To remove the navigation symbols from the bottom of all slides

\graphicspath{{../img/}}
\usepackage[utf8]{inputenc}
\usepackage[T1]{fontenc} %% https://tex.stackexchange.com/questions/664/why-should-i-use-usepackaget1fontenc
\usepackage{graphicx}
\usepackage{listings}
\usepackage{colortbl}
\usepackage{tikz}
\usetikzlibrary{trees}

\title{¿Cómo contribuir al kernel de Linux?}
\author[\textcircled{cc} BY-SA]{Miguel Bernal Marin} % Your name

\institute[]
{
Cumbre de Contribuidores de Open Source Software  \\
https://github.com/miguelinux/ccoss
}
\date{
    20 de Septiembre de 2022
}

\logo{\includegraphics[height=1cm]{contribu-logo.png}}
\newcommand{\nologo}{\setbeamertemplate{logo}{}} % command to set the logo to nothing

%New colors defined below
\definecolor{codeNumbers}{rgb}{0.5,0.5,0.5}
\definecolor{editorGray}{rgb}{0.95, 0.95, 0.95}
\definecolor{editorOcher}{rgb}{1, 0.5, 0} % #FF7F00 -> rgb(239, 169, 0)
\definecolor{editorGreen}{rgb}{0, 0.5, 0} % #007C00 -> rgb(0, 124, 0)
\definecolor{brown}{rgb}{0.69,0.31,0.31}

%Code listing style named "mystyle"
\lstdefinestyle{mystyle}{
  % General design
  backgroundcolor=\color{editorGray},
  basicstyle={\ttfamily},
  numberstyle=\tiny\color{codeNumbers},
  xleftmargin={0.2cm},
  numbersep=5pt,
  numbers=left,
  stepnumber=1,
  firstnumber=1,
  numberfirstline=true,
  % Code design
  identifierstyle=\color{black},
  keywordstyle=\color{blue}\bfseries,
  ndkeywordstyle=\color{editorGreen}\bfseries,
  stringstyle=\color{editorOcher}\ttfamily,
  commentstyle=\color{brown}\ttfamily,
  % Code
  alsodigit={.:;},
  tabsize=2,
  showtabs=false,
  showspaces=false,
  showstringspaces=false,
  extendedchars=true,
  breaklines=true,
  % Español
  literate=%
  {á}{{\'a}}1 {é}{{\'e}}1 {í}{{\'i}}1 {ó}{{\'o}}1 {ú}{{\'u}}1
  {Á}{{\'A}}1 {É}{{\'E}}1 {Í}{{\'I}}1 {Ó}{{\'O}}1 {Ú}{{\'U}}1
  {¡}{{!`}}1  {¿}{{?`}}1
  {ñ}{{\~n}}1 {Ñ}{{\~N}}1
}

%"mystyle" code listing set
\lstset{style=mystyle}

%------------------------------------------------------------
%The next block of commands puts the table of contents at the
%beginning of each section and highlights the current section:
\AtBeginSection[]
{
  \begin{frame}
    \frametitle{Contenido}
    \tableofcontents[currentsection]
  \end{frame}
}
%------------------------------------------------------------

\newcommand{\pausa}{\pause} % Para usar una pausa en las presentaciones
%\newcommand{\pausa}{}      % Para que NO salgan las pausas

\begin{document}

\begin{frame}
    \titlepage
\end{frame}

{\nologo% ex: ts=2 sw=2 sts=2 et filetype=tex
% SPDX-License-Identifier: CC-BY-SA-4.0

\begin{frame}[c]{Legal Disclaimer}
  This presentation is for informational purposes only.
  INTEL MAKES NO WARRANTIES, EXPRESS OR IMPLIED, IN THIS SUMMARY.

  \vspace{\baselineskip}
  Intel and the Intel logo are trademarks of Intel Corporation in the
  U.S. and/or other countries.

  \vspace{\baselineskip}
  * Other names and brands an logos may be claimed as the property of others.

  \vspace{\baselineskip}
  Copyright © 2021, Intel Corporation. All rights reserved.
\end{frame}
}
% ex: ts=2 sw=2 sts=2 et filetype=tex
% SPDX-License-Identifier: CC-BY-SA-4.0

\begin{frame}[c]{¡Hola!}

  \begin{columns}
    \column{0.7\textwidth}
      Miguel Bernal Marin

      \vspace{\baselineskip}
      OS Developer Engineer \textbf{@Intel}
    \column{0.3\textwidth}
      \includegraphics[scale=0.2]{miguelinux-color.png}
  \end{columns}

  Profesor:
  \begin{itemize}
    \item \href{http://tecmm.edu.mx/}{Instituto Tecnológico Mario
      Molina Pasquel y Henriquez (TecMM)}
    \item \href{https://iteso.mx/}{ITESO}
  \end{itemize}

  \begin{center}
    \href{http://tecmm.edu.mx/}{\includegraphics[scale=0.5]{tecmm-z-logo.png}}
    \hspace{2cm}
    \href{https://iteso.mx/}{\includegraphics[scale=1.8]{iteso-logo.png}}
  \end{center}
\end{frame}


\begin{frame}
    \frametitle{Contenido}
    \tableofcontents
\end{frame}

% ex: ts=2 sw=2 sts=2 et filetype=tex
% SPDX-License-Identifier: CC-BY-SA-4.0

\section{¿Qué es Linux?}

\begin{frame}[c]{Historia}
  \begin{columns}
    \column{0.6\textwidth}
     \begin{itemize}
      \item Linux comienza en 1991 como un proyecto personal del estudiante
        finlandés \textbf{Linus Benedict Torvalds}$^{\hyperlink{creditos}{[1]}}$.
      \pausa
      \item Crea un nuevo núcleo de un sistema operativo libre, bajo la
        \textbf{GNU General Public License v2} (GPL v2)
      \pausa
      \item El núcleo Linux resultante ha estado marcado por un crecimiento
        constante a lo largo de su historia.
     \end{itemize}
    \column{0.4\textwidth}
      \pausa
      \begin{center}
        \includegraphics[scale=0.6]{lfs/linus.png}
      \end{center}
  \end{columns}
\end{frame}

\begin{frame}[fragile]
  \frametitle{El segundo anuncio en usenet$^{\hyperlink{creditos}{[2]}}$}
  \begin{block}{}
  \begin{verbatim}
Newsgroups: comp.os.minix
Subject: What would you like to see most in minix?
Date: 25 Aug 91 20:57:08 GMT

Hello everybody out there using minix -

I'm doing a (free) operating system (just a hobby, won't be big and
professional like gnu) for 386(486) AT clones.  This has been brewing
since april, and is starting to get ready.  I'd like any feedback on
things people like/dislike in minix, as my OS resembles it somewhat
(same physical layout of the file-system (due to practical reasons)
among other things).
  \end{verbatim}
  \end{block}
\end{frame}

\begin{frame}[c]{Historia de las versiones del kernel$^{\hyperlink{creditos}{[3]}}$}
  \begin{center}
    \includegraphics[scale=0.39]{lfs/linux_kernel_timeline.png}
  \end{center}
\end{frame}

\begin{frame}[c]{Versiones activas del kernel}
  Hay varias categorías principales en las que pueden caer los
  lanzamientos del kernel$^{\hyperlink{creditos}{[4]}}$:

  \vspace{\baselineskip}
  \begin{description}
    \item[Prepatch] Los kernels de preparche o "RC" (candidanto a
      liberación) son versiones preliminares del kernel de la línea principal
      (mainline) que están dirigidas principalmente a otros desarrolladores
      de kernel y entusiastas de Linux.

      \vspace{\baselineskip}
      Deben compilarse desde la fuente y, por lo general, contienen nuevas
      características que deben probarse antes de que puedan colocarse en
      una versión estable.

      \vspace{\baselineskip}
      Linus Torvalds mantiene y publica los kernels previos al parche.
  \end{description}
\end{frame}

\begin{frame}[c]{Versiones activas del kernel}
  \begin{description}
    \item[Mainline] Linus Torvalds mantiene el árbol principal (mainline).
      Es el árbol donde se presentan todas las funciones nuevas y donde
      ocurre todo el desarrollo nuevo y emocionante.

      \vspace{\baselineskip}
      Los nuevos núcleos principales se lanzan cada 9 o 10 semanas.
  \end{description}
\end{frame}

\begin{frame}[c]{Versiones activas del kernel}
  \begin{description}
    \item[Stable] Después de que se lanza cada núcleo principal (mainline),
      se considera "estable".

      \vspace{\baselineskip}
      Todas las correcciones de errores para un kernel estable se
      retroalimentan (backported) desde el árbol de la línea principal
      y las aplica un mantenedor de kernel estable designado.

      \vspace{\baselineskip}
      Por lo general, solo hay unas pocas versiones del kernel con corrección
      de errores hasta que el siguiente kernel de la línea principal esté
      disponible, a menos que se designe como un
      "kernel de mantenimiento a largo plazo".

      \vspace{\baselineskip}
      Las actualizaciones estables del kernel se publican según sea
      necesario, generalmente una vez a la semana.
  \end{description}
\end{frame}

\begin{frame}[c]{Versiones activas del kernel}
  \begin{description}
    \item[Longterm] Por lo general, se proporcionan varias versiones del
      kernel de "\textbf{mantenimiento a largo plazo}" con el fin de
      respaldar las correcciones de errores para los árboles del kernel
      más antiguos.

      \vspace{\baselineskip}
      Solo se aplican correcciones de errores importantes a dichos núcleos
      y, por lo general, no se lanzan con mucha frecuencia,
      especialmente para los árboles más antiguos.
  \end{description}
\end{frame}

\begin{frame}[c]{Núcleos de lanzamiento a largo plazo}
  \begin{table}[]
  \begin{tabular}{llll}
    \textbf{Versión} &  \textbf{Mantenedor} & \textbf{Liberado} & \textbf{Fin de vida} \\
     & & & \textbf{proyectado} \\
    \rowcolor{editorGray}
    5.15&Greg Kroah-Hartman \& Sasha Levin 	&2021-10-31 	&Oct, 2023 \\
    5.10&Greg Kroah-Hartman \& Sasha Levin 	&2020-12-13 	&Dec, 2026 \\
    \rowcolor{editorGray}
    5.4&Greg Kroah-Hartman \& Sasha Levin 	&2019-11-24 	&Dec, 2025 \\
    4.19&Greg Kroah-Hartman \& Sasha Levin 	&2018-10-22 	&Dec, 2024 \\
    \rowcolor{editorGray}
    4.14&Greg Kroah-Hartman \& Sasha Levin 	&2017-11-12 	&Jan, 2024 \\
    4.9&Greg Kroah-Hartman \& Sasha Levin 	&2016-12-11 	&Jan, 2023 \\
  \end{tabular}
  \end{table}
\end{frame}



\begin{frame}[c]{Historia de las versiones del kernel 4.x$^{\hyperlink{creditos}{[3]}}$}
  \begin{center}
    \includegraphics[scale=0.7]{lfs/linux_kernel_timeline_4.png}
  \end{center}
\end{frame}

\begin{frame}[c]{Historia de las versiones del kernel 5.x$^{\hyperlink{creditos}{[3]}}$}
  \begin{center}
    \includegraphics[scale=0.7]{lfs/linux_kernel_timeline_5.png}
  \end{center}
\end{frame}


%\section{}
%
%\begin{frame}[c]{}
%  \vspace{\baselineskip}
%\end{frame}
%
%\begin{frame}[fragile]
%  \frametitle{ufw - Cortafuegos sin complicaciones}
%  \begin{lstlisting}[language=Bash,numbers=none]
%  \end{lstlisting}
%\end{frame}


\begin{frame}[c]{Créditos}
  \label{creditos}
  El material está basado en
  \begin{enumerate}
    \item
      \href{https://es.wikipedia.org/wiki/Historia_de_Linux}
      {https://es.wikipedia.org/wiki/Historia\_de\_Linux}
    \item
      \href{https://www.cs.cmu.edu/~awb/linux.history.html}
      {https://www.cs.cmu.edu/\~{}awb/linux.history.html}
    \item
      \href{https://en.wikipedia.org/wiki/Linux_kernel_version_history}
      {https://en.wikipedia.org/wiki/Linux\_kernel\_version\_history}
    \item
      \href{https://kernel.org/category/releases.html}
      {https://kernel.org/category/releases.html}
    \item
      \href{https://www.kernel.org/doc/html/latest/process/index.html}
      {https://www.kernel.org/doc/html/latest/process/index.html}
    \item
      \href{https://www.kernel.org/cloning-linux-from-a-bundle.html}
      {https://www.kernel.org/cloning-linux-from-a-bundle.html}
    \item
      \href{https://www.kernel.org/doc/html/latest/process/howto.html}
      {https://www.kernel.org/doc/html/latest/process/howto.html}
    \item
      \href{https://www.kernel.org/doc/html/latest/admin-guide/reporting-issues.html}
      {https://www.kernel.org/doc/html/latest/admin-guide/reporting-issues.html}
%    \item
%      \href{}
%      {}
  \end{enumerate}
\end{frame}

\end{document}
